%% ============================================================
%%  Chapitre 5 — Architecture multi-agents
%% ============================================================
\chapter{Architecture multi-agents}
\label{ch:agents}

\section{Principes de conception}

L'architecture multi-agents de \textit{Scientific Watch Agent} repose sur
quatre principes directeurs.

\paragraph{Séparation des responsabilités.}
Chaque agent a une responsabilité unique et bien délimitée : recherche,
filtrage, résumé, synthèse, critique, analyse. Cette séparation facilite
le test, la maintenance et l'évolution indépendante de chaque composant.

\paragraph{Contrat de communication formalisé.}
Tous les échanges entre agents transitent par un état partagé (\texttt{WatchState})
dont chaque champ est typé par un schéma Pydantic v2. Ce contrat formalisé
garantit la cohérence des données tout au long du pipeline et constitue
une documentation vivante de l'interface entre agents.

\paragraph{Principe Single Writer / Multiple Readers.}
Chaque section de l'état partagé n'est écrite que par un unique agent,
mais peut être lue par tous. Cette règle prévient les conditions de course
(\textit{race conditions}) dans les architectures futures à agents concurrents.

\paragraph{Abstraction du fournisseur LLM.}
Les agents ne connaissent pas le fournisseur de LLM qu'ils utilisent.
Ils invoquent une interface unifiée \texttt{LLMClient} qui délègue
automatiquement au fournisseur approprié selon la tâche, via une fabrique
(\texttt{factory.py}).

\section{État partagé : \texttt{WatchState}}

L'état partagé \texttt{WatchState} est défini comme un \texttt{TypedDict}
Python dans \texttt{src/agents/state.py}. Il contient :

\begin{itemize}
  \item \texttt{topic} : le sujet de recherche soumis par l'utilisateur ;
  \item \texttt{expanded\_queries} : liste des requêtes générées par
    le \texttt{QueryExpander} ;
  \item \texttt{raw\_articles} : liste d'objets \texttt{Article} récupérés
    par le \texttt{Searcher} ;
  \item \texttt{scored\_articles} : liste d'articles avec leur
    \texttt{QualityScore} associé, produite par le \texttt{QualityCritic} ;
  \item \texttt{summaries} : liste d'\texttt{ArticleSummary} ou
    \texttt{NarrativeSummary} générés par le \texttt{Summarizer} ;
  \item \texttt{synthesis} : objet \texttt{Synthesis} produit par
    le \texttt{Synthesizer} ;
  \item \texttt{critic\_feedback} : objet \texttt{CriticFeedback}
    du dernier cycle Reflexion ;
  \item \texttt{reflexion\_count} : compteur d'itérations Reflexion ;
  \item \texttt{trend\_analysis} : objet \texttt{TrendAnalysis} final ;
  \item \texttt{logs}, \texttt{errors} : listes accumulées via
    \texttt{Annotated[list, operator.add]}.
\end{itemize}

\section{Couche d'abstraction LLM}

\subsection{Interface \texttt{LLMClient}}

L'interface \texttt{LLMClient} est définie comme un Protocol Python
dans \texttt{src/llm/base.py}. Elle expose deux méthodes :

\begin{lstlisting}[caption={Interface LLMClient (src/llm/base.py).}]
class LLMClient(Protocol):
    def chat(
        self,
        system: str,
        messages: list[dict],
        temperature: float = 0.3,
        max_tokens: int = 2048,
    ) -> LLMResponse: ...

    def chat_structured(
        self,
        system: str,
        messages: list[dict],
        schema: type[T],
        temperature: float = 0.1,
    ) -> tuple[T, LLMResponse]: ...
\end{lstlisting}

\subsection{Quatre fournisseurs implémentés}

Quatre implémentations de \texttt{LLMClient} sont disponibles,
présentées dans le tableau~\ref{tab:providers}.

\begin{table}[htbp]
  \centering
  \caption{Fournisseurs LLM intégrés et méthode de structured output.}
  \label{tab:providers}
  \rowcolors{2}{lightgray}{white}
  \begin{tabular}{llll}
    \toprule
    \textbf{Provider} & \textbf{Modèles} & \textbf{Structured output} & \textbf{Coût relatif} \\
    \midrule
    Anthropic & claude-sonnet-4-6 & \texttt{tool\_use} forcé & $\times 10$ \\
    OpenAI & gpt-4o-mini & \texttt{response\_format} JSON & $\times 3$ \\
    Groq & llama-3.3-70b & \texttt{response\_format} JSON & $\times 1$ \\
    Nebius & DeepSeek-V3.2, Qwen3 & \texttt{response\_format} JSON & $\times 0.5$ \\
    \bottomrule
  \end{tabular}
\end{table}

La fabrique \texttt{get\_llm\_for\_task(task)} dans \texttt{src/llm/factory.py}
sélectionne automatiquement le modèle adapté à chaque type de tâche
selon la table \texttt{TASK\_MODELS} définie dans \texttt{src/config.py}.
Cette approche permet d'optimiser le rapport coût/qualité : Claude Sonnet
pour la synthèse (contexte 200K), Nebius Qwen3 pour la summarisation
répétitive (\$0.14/M tokens), GPT-4o-mini pour l'expansion de requêtes.

\section{Les sept agents}
\label{sec:agents}

\subsection{QueryExpander — Expansion de requêtes (ReAct)}

Le \texttt{QueryExpander} est le premier agent du pipeline. Il reçoit
le topic brut de l'utilisateur et génère plusieurs requêtes de recherche
diversifiées, optimisées pour maximiser la couverture des travaux pertinents.
Il implémente le pattern ReAct (détaillé au Chapitre~\ref{ch:patterns}).

À chaque itération, l'agent :
\begin{enumerate}
  \item Raisonne sur la meilleure requête à formuler ;
  \item Sonde OpenAlex avec cette requête ;
  \item Observe le nombre d'articles retournés et les concepts associés ;
  \item Adapte sa stratégie (affiner, élargir ou diversifier).
\end{enumerate}

À l'issue de 4 itérations maximum, le \texttt{QueryExpander} produit
une liste de 3 à 8 requêtes complémentaires qui alimentent le \texttt{Searcher}.

\subsection{Searcher — Récupération et déduplication}

Le \texttt{Searcher} exécute en parallèle toutes les requêtes produites
par le \texttt{QueryExpander} sur l'API OpenAlex, fusionne les résultats
et déduplique par identifiant OpenAlex. Il enrichit également les données
auteurs (h-index, affiliations) via des appels secondaires à l'API.

Le corpus résultant contient typiquement entre 80 et 300 articles, selon
la richesse du domaine et le nombre de requêtes générées.

\subsection{QualityCritic — Filtrage par qualité}

Le \texttt{QualityCritic} applique le scorer multi-dimensionnel
(Chapitre~\ref{ch:automl}) à chaque article du corpus. Il charge
automatiquement les pondérations Optuna apprises pour le domaine correspondant
(via \texttt{load\_weights\_for\_topic}), ou utilise les pondérations par
défaut si aucune pondération spécifique n'est disponible.

Il retient les $k$ articles les mieux classés (typiquement $k = 15$)
pour alimenter le \texttt{Summarizer}.

\subsection{Summarizer — Résumés structurés et narratifs}

Le \texttt{Summarizer} génère, pour chaque article sélectionné, un résumé
dans l'un des deux modes suivants :

\paragraph{Mode structuré.}
Le résumé est un objet \texttt{ArticleSummary} comportant six champs :
\textit{problem} (problème adressé), \textit{method} (approche proposée),
\textit{dataset} (données utilisées), \textit{results} (résultats obtenus),
\textit{limitations} (limites identifiées), et \textit{contribution}
(apport principal).

\paragraph{Mode narratif.}
Le résumé est un paragraphe en prose de 150 à 250 mots (\texttt{NarrativeSummary}),
plus accessible et adapté à une présentation à des lecteurs non spécialistes.

\subsection{Synthesizer — Synthèse globale (Reflexion)}

Le \texttt{Synthesizer} reçoit l'ensemble des résumés et génère une
synthèse globale du domaine (\texttt{Synthesis}) comprenant : une vue
d'ensemble du domaine, les principales approches méthodologiques,
les datasets de référence, et les résultats clés.

Il est au cœur du pattern Reflexion (Chapitre~\ref{ch:patterns}) :
sa production peut être révisée itérativement en réponse aux retours du
\texttt{Critic}.

\subsection{Critic — Évaluation Reflexion}

Le \texttt{Critic} évalue la synthèse produite par le \texttt{Synthesizer}
selon quatre axes : \textit{fidelity} (fidélité aux sources), \textit{completeness}
(exhaustivité), \textit{specificity} (niveau de détail), et \textit{consistency}
(cohérence interne). Il produit un objet \texttt{CriticFeedback} indiquant
si une révision est nécessaire (\texttt{needs\_revision}) et, le cas échéant,
les améliorations à apporter.

\subsection{TrendAnalyst — Analyse des tendances}

Le \texttt{TrendAnalyst} reçoit la synthèse finale et génère un objet
\texttt{TrendAnalysis} comprenant : les tendances émergentes du domaine
(avec leur niveau de maturité), les lacunes de recherche identifiées,
et des perspectives de directions futures. Un mécanisme de retry automatique
garantit que le champ \texttt{future\_perspectives} n'est jamais vide.

\section{Flux d'exécution et arêtes conditionnelles}

Le pipeline est compilé par LangGraph en un graphe orienté acyclique
(DAG) avec une arête conditionnelle pour le cycle Reflexion.
La figure~\ref{fig:langgraph_pipeline} illustre le flux complet.

\begin{figure}[htbp]
  \centering
  \includegraphics[width=0.85\textwidth]{figures/langgraph_pipeline.pdf}
  \caption{Pipeline LangGraph avec ses sept agents et l'arête conditionnelle
  Reflexion. L'arête \texttt{route\_after\_critic} redirige vers le
  \texttt{Synthesizer} si \texttt{needs\_revision = True} et que le
  nombre d'itérations est inférieur au maximum autorisé (3).}
  \label{fig:langgraph_pipeline}
\end{figure}

La fonction \texttt{route\_after\_critic} dans \texttt{src/agents/graph.py}
implémente la logique de routage :
\begin{itemize}
  \item Si \texttt{critic\_feedback.needs\_revision = True} et
    \texttt{reflexion\_count < MAX\_REFLEXION\_ITERATIONS} : retour au
    \texttt{Synthesizer} ;
  \item Sinon (qualité suffisante ou itérations épuisées) : passage au
    \texttt{TrendAnalyst}.
\end{itemize}

\section{Évaluation du pipeline}

Un exemple d'exécution du pipeline complet sur le topic
``\textit{fake news detection}'' produit les métriques suivantes
(tableau~\ref{tab:pipeline_demo}) :

\begin{table}[htbp]
  \centering
  \caption{Métriques d'exécution typiques du pipeline (topic: fake news detection).}
  \label{tab:pipeline_demo}
  \rowcolors{2}{lightgray}{white}
  \begin{tabular}{lrr}
    \toprule
    \textbf{Agent} & \textbf{Durée (s)} & \textbf{Coût (\$)} \\
    \midrule
    QueryExpander (4 iter.) & 12.3 & 0.002 \\
    Searcher (3 requêtes) & 8.7 & 0.000 \\
    QualityCritic (120 articles) & 0.3 & 0.000 \\
    Summarizer (15 articles) & 45.2 & 0.031 \\
    Synthesizer + Critic (2 iter.) & 18.9 & 0.012 \\
    TrendAnalyst & 7.4 & 0.005 \\
    \midrule
    \textbf{Total} & \textbf{92.8} & \textbf{0.050} \\
    \bottomrule
  \end{tabular}
\end{table}

Le coût total d'un run est de l'ordre de \$0.05 avec la configuration
Nebius/Groq, ce qui est compatible avec un usage régulier dans le cadre
d'un projet de recherche académique.
